\chapter{Scarcity, access and unequal exposure}\label{ch:scarcity}

\section{A full store and an empty cup}
At Ridge, the storekeeper can see water through the inspection window. At Vale, a resident opens a tap and receives none. Before deciding that the valley has ``run out of water'', we need to identify what has failed. The stock may be insufficient, the pipe may be broken, the water may be unsuitable, or the resident may be unable to obtain permission or afford delivery. Each situation can produce the same immediate experience: an empty cup.

The distinction changes the response. Repairing a pipe does not solve the absence of water at its source. Increasing the source stock does not fix a locked gate. Reducing a delivery charge does not remove contamination. A diagnosis expressed only as scarcity is incomplete until the relevant resource, location, time, quality and access arrangement are specified.

FAO's water-scarcity overview explicitly distinguishes physical scarcity, access scarcity and infrastructure scarcity. Its categories separate insufficient suitable water from water that exists but is inaccessible because of institutional arrangements or missing delivery infrastructure. The distinction gives a useful real-world reason to avoid treating a regional total as a complete account of provision.\cite{intro-water}

Our invented valley will let us practise that distinction without pretending that a two-place model represents every cause of deprivation. In particular, a stock variable does not by itself encode a person's income, legal status, disability, available time or ability to travel. Those matters may determine access even when the physical stock is measured correctly.

\section{Scarcity relative to a task}
The statement ``there are six units'' is a quantity statement. The statement ``six units are enough'' relates that quantity to a purpose. Six may be enough for one declared task and insufficient for another. It may be enough now but insufficient over a longer interval. It may be available in aggregate but divided among places in a way that makes the intended service impossible.

This is why need and demand require careful language. A \emph{need} identifies something required for a declared human or functional purpose. A \emph{request} expresses a desired service to a provider. An \emph{effective market demand} additionally involves willingness and ability to transact under particular terms. These concepts may coincide in an example, but they are not definitions of one another.

Suppose a clinic requests four units and a decorative fountain requests four units. Equal quantities do not establish equal urgency. Nor does a larger offer of payment prove a larger health consequence if the request is refused. A priority rule might account for purpose and urgency, but that rule must be declared. It cannot be extracted from the arithmetic fact that four equals four.

The physical teaching model can represent quantities and transfers faithfully while leaving such priorities open. This may feel unsatisfying at first. Yet silently embedding a priority rule would be less satisfactory, because it would prevent readers from seeing where a decision entered. A transparent research programme should state both the represented physical comparison and the additional institutional choices needed to use it.

\section{Stocks, routes and access form different bottlenecks}
Consider three hypothetical mornings. On the first, Ridge holds too little suitable water to meet all declared uses. On the second, Ridge holds enough, but the only pipe is closed for repair. On the third, the pipe works and the stock is sufficient, but a household cannot obtain delivery under the payment rule. In each case, adding the same extra quantity at Ridge may have a different effect.

On the first morning, added suitable supply could directly relieve the physical shortfall, subject to distribution. On the second, it may merely enlarge an already adequate inaccessible store. On the third, it may leave the household excluded. These are logical consequences of the declared situations; they are not measured estimates of how often each bottleneck occurs.

A useful account can therefore include a chain of questions. Is the resource present? Is it suitable? Is a route available? Can the route carry the required quantity in time? Is someone authorised and able to perform the action? Can the intended person use the resulting service? The first negative answer may explain the immediate failure, but repairing it can reveal another constraint farther along the chain.

This chain helps explain why a scalar score is never a substitute for inspection. If the score improves because more resource is stored upstream while the delivery route remains broken, the relevant function may not improve. A model intended to evaluate that function would need to represent the relationship adequately. A limited stock model should describe its conclusion as limited stock redistribution.

\section{Unequal exposure means more than unequal totals}
Two households may face the same service interruption and experience different consequences. One may have stored supplies, flexible work and transport. Another may have none of these and include a person whose care depends on reliable access. Equality of the external interruption does not imply equality of exposure, coping capacity or loss.

The World Health Organization's 2025 report on social determinants of health equity identifies living conditions and access to power, money and resources as important influences on avoidable health gaps. Its overview describes variation associated with social position and circumstances; it does not reduce every health outcome to income alone.\cite{intro-health}

For the present book, the lesson is that socioeconomic conditions should not disappear behind a regional average. A hypothetical average delivery of ten units can arise when everyone receives ten or when one group receives twenty and another receives zero. Those are the same average and different arrangements. Whether they are acceptable depends on purpose, needs and rights as well as arithmetic.

The phrase ``socioeconomic inequality'' also covers several dimensions. Income concerns a flow over time; wealth concerns accumulated claims or assets; access concerns practical ability to obtain a service; power concerns influence over rules and decisions. Combining them into one vague disadvantage measure may obscure how a particular failure can be repaired.

EBU's physical account cannot settle all these dimensions. Indeed, an institution that rewarded only easily measured actions could create new exclusions. Its design would need to ask who can participate, whose work is observable, who bears measurement costs and who can challenge an incorrect record. These questions arise from the proposal's intended use; they are not answered by a favourable conservation residual.

\section{A total can conceal a distribution}
Return to our fixed twenty-unit example. Ridge and Vale can hold $(18,2)$, $(14,6)$ or $(10,10)$. The total remains twenty. We can write
\begin{equation}
18+2=14+6=10+10=20.
\end{equation}
The equality says that the represented quantity is conserved. It does not say that the arrangements provide the same service. A particular reference may prefer one distribution, but the reference must be stated rather than inferred from conservation.

Suppose, for illustration, that ten units at each place are the declared comparison configuration. Then $(10,10)$ matches that declaration and $(18,2)$ does not. Now suppose Ridge houses a much larger function than Vale and the declared compatible reference is instead $(14,6)$. The same physical arrangements receive a different interpretation. We have changed the evaluation, not the total mass.

Neither reference is justified by being easy to write. An applied reference would require reasons connected to the represented function. It might need evidence about service requirements, physical capacity or agreed provision. The introductory arithmetic does not supply such evidence. Its purpose is to make the point of choice visible before the later potential assigns a smooth numerical burden.

\begin{intuitionbox}{Distribution is visible only when places remain visible}
An aggregate can be conserved while one place is deprived. Retaining separate state coordinates allows us to examine distribution. Choosing how to evaluate those coordinates requires a declared reference and does not follow from the conservation law.
\end{intuitionbox}

\section{The cost of reaching a service}
Access is not always a yes-or-no gate at a shop. Consider a hypothetical patient who can obtain medicine only by making a long journey. The medicine is present, the price is affordable, and the transaction is legal. Yet the journey may require time away from paid work, assistance from another person and reliable transport. A simple inventory reports availability while the patient faces a substantial practical obstacle.

This example shows why the boundary of an account matters. If the account stops at the warehouse, it can correctly record medicine in stock while missing the final access problem. If it extends to the clinic, it can record arrival there while still missing the patient's journey. Every extension must be justified by the question being studied and supported by relevant data.

There is no honest shortcut in which the word ``local'' makes all these consequences disappear. Local observation can be sufficient for a specified mathematical field difference when all affected terms are known. It is not a guarantee that the representation includes every socially important effect. The local theorem and the adequacy of the model are different questions.

Nor should the patient be charged an invented physical burden simply because several people helped. A route has real transformations and real resource uses, each requiring its own declared account. Counting the same electricity once at generation, again at delivery and again as an unexplained surcharge would exaggerate burden. Linked actions should reveal a chain, not multiply one effect by the number of records.

\section{Essential support and the danger of circular qualification}
Imagine an institution that says a person may receive water only after generating enough positive EBU. A person who is ill and unable to act could then lose the support needed to recover. The qualification rule would create a circle: access requires contribution, while contribution requires access. No physical accounting identity resolves that institutional choice.

The ethical position adopted in this introduction is that essential support deserves explicit protection. The precise mechanism remains a matter for institutional design and public justification. The important mathematical boundary is that an action's signed field result is not a measure of a person's worth or a complete entitlement rule.

Similarly, a settlement may need to make a physically costly action available because the alternative is unacceptable. Delivering medicine over a difficult route may increase some represented burdens while meeting an urgent need. The account should preserve that information. Hiding the burden to make the action look favourable would obstruct learning; denying the need because the burden is unfavourable would confuse evaluation with permission.

The better question is what the account helps the community learn. Perhaps a local store would reduce repeated emergency journeys. Perhaps a more reliable route would help. Perhaps the chosen potential omits a crucial function. These possibilities require investigation. They do not follow automatically from the sign of a single action.

\section{What a fair comparison would have to show}
If EBU is eventually tested as guidance for a resource arrangement, an apparent improvement in a total cannot be the whole result. A comparison needs to report who received service, which constraints were respected and what happened under difficult conditions. It also needs to examine cases in which EBU guidance performs poorly.

For example, a rule that never transfers anything can preserve every source stock over a short action-only interval. That does not establish satisfactory service. A rule that serves every immediate request may deplete a source needed elsewhere. A rule that improves an aggregate potential may leave a particular essential request unmet. Each behaviour reveals a different dimension.

The book does not select a final controller or freeze study metrics here. Instead, it develops the account needed to state such questions precisely. A research design should declare its service, physical and distributional criteria before observing outcomes. Otherwise, the evaluator can choose whichever interpretation makes its favoured rule look successful.

\section{The deadline changes the access question}
Two deliveries bring the same quantity to Vale. One arrives before the clinic opens; the other arrives after the day's procedures have been cancelled. An end-of-day stock report may show the same quantity in both cases. The service histories are different. Access therefore includes timing whenever the function has a deadline.

To represent that difference, an account needs an appropriate time resolution and a declared service record. It would be misleading to repair the problem simply by assigning a more attractive final reference after the late delivery. The reference evaluates a state; a deadline concerns whether a service occurred during a specified interval. Both may matter, but neither substitutes for the other.

Timing can also interact with unequal circumstances. In an invented example, one resident can return tomorrow, while another has only one available journey. A uniform postponement can therefore impose different practical burdens. Establishing those burdens in a real setting requires evidence about the affected people, rather than a presumption based on their group label.

This is a useful design lesson for future studies. If sustained demand is the intended subject, the study must preserve the time structure of demand and delivery. A total delivered eventually may be an inadequate outcome. The criterion should be declared before results are inspected, so that delayed service is neither excused nor condemned according to which controller produced it.

\section{Common misunderstandings}
``Scarcity is always just a shortage of total material'' overlooks location, timing, suitability and access. Conversely, ``all scarcity is artificial'' overlooks genuine physical limits. A careful diagnosis can recognise both without treating either as the explanation of every case.

``Equal distribution is automatically fair'' mistakes one possible numerical pattern for a complete ethical argument. Equal quantities may support unequal functions or needs. A reference must be justified in its intended setting, and a scalar allocation cannot settle every dimension of fairness.

``A positive physical score proves that inequality has improved'' makes a claim the score may not represent. If access or distributional outcomes matter, the account must measure them separately or explicitly include a justified representation. A label cannot add missing information.

``The model is useless unless it includes everything'' sets an impossible standard. A small model can expose an accounting error or prove a conditional identity. Its usefulness depends on maintaining the connection between the narrow result and the question it actually answers.

\section{Guided problems and answers}
\subsection{Three identical complaints}
Three residents say, ``There is no water.'' In one case the source is empty, in another the pipe is broken, and in the third delivery is refused for nonpayment. Would the same intervention solve all three?

Not necessarily. The first requires suitable supply or a change in uses; the second requires a route or repair; the third requires an access decision under the relevant institution. Each intervention can also encounter additional constraints. A useful record identifies the failure rather than replacing all three with a common headline.

\subsection{The average service report}
Two places each require attention. A report states that average delivery was ten units. Can we infer that each received ten?

No. Many pairs have that average, including $(20,0)$ and $(10,10)$. The distribution must be reported if it matters to the claim. We must also establish whether equal delivery is the relevant goal. Averages summarise information; they do not preserve every fact from which they were calculated.

\subsection{The unfavourable essential action}
An accepted emergency delivery has a negative signed EBU result in the declared model. Must the delivery be refused?

The sign alone supplies no such rule. Permission, essential support and institutional financing need their own declarations. The result should remain visible so that future arrangements can be improved and the model's adequacy questioned. A truthful account can coexist with an explicit decision to perform a costly necessary action.

\section{What to carry forward}
Scarcity can concern physical quantity, infrastructure or access, and unequal circumstances can turn the same interruption into very different experiences. Conservation is essential, but totals alone cannot establish provision. These distinctions prepare us to consider planetary pressures without losing sight of the people and functions through which those pressures are experienced.
